% P10-4c, 14.09.2026: Ledger-Herleitung und Nachweisanschlüsse präzisiert.
% M07, 11.09.2026: Begriffe vor Verarbeitungstabelle erklärt;
% Adjudikationsvorlage und datierten Facettennachtrag aus 5.6 hierher
% verschoben. Selektive Evidenzklärung bleibt von Autorisierung getrennt.
% Belege: SemanticLedgerEntry / ArtifactDisposition;
% EvidenceLedgerProjection.cs; F1 consumable.json; rekonstruierte Queue
% thesis-evidence/w2/hitl-queue-8189ea/queue.md, US::AU-0124;
% LedgerAdjudicationReviewAdapter.UsActionOptions;
% Facetten-Integrationslauf 20260909_190518_f58d6a (D.6).

\section{Beobachtbare und evidenzgebundene Informationsverarbeitung}
\label{sec:evidenzgebundene-verarbeitung}

Die erste Säule bündelt die Anforderungen A1 bis A3. Der
\emph{Evidence Ledger} konkretisiert für transkriptbasierte Eingaben
die in Abschnitt~\ref{sec:evidenzgebundene-wissensgewinnung}
begründete Quellenstrukturierung: Zwischen Rohtranskript und
fachlicher Ableitung hält er Aussagen und ihre Herkunft in einer
überprüfbaren Zwischenrepräsentation fest. Er entspricht der
Evidenzgewinnung des Meeting-Pfads aus
Abschnitt~\ref{sec:logische-gesamtarchitektur}. Das Rohtranskript
bleibt die ursprüngliche Quelle; der maßgebliche fachliche
Projektzustand wird getrennt davon im Core geführt.

\emph{Evidence Units} sind einzeln adressierbare Quellausschnitte,
die durch deterministische Segmentierung des Transkripts entstehen.
Aus ihnen extrahiert das Modell fachlich relevante Aussagen, die
\emph{Claims}. Jeder Claim besitzt einen Aussagentext
(\emph{Proposition}), eine Kennung sowie Bezüge zu den
zugrunde liegenden Units und gespeicherte Quellzitate. Die in
Abbildung~\ref{fig:evidenz-provenienzmodell} dargestellte
Claim-Kennung ermöglicht Referenzen innerhalb des betreffenden Laufbestands;
die unterschiedlichen Kennungsreichweiten erläutert
Abschnitt~\ref{sec:core-projektfortschreibung}.
Eine anschließende modellgestützte \emph{Kanonisierung} führt die
Kandidaten in einen konsolidierten Claim-Bestand über und kann eng
zusammengehörige Inhalte bündeln. Ein Claim ist daher nicht zwingend
mit einer atomaren Referenzaussage der Evaluation identisch.

\emph{Facetten} beschreiben Eigenschaften einer Aussage, etwa ihren
fachlichen Typ, ihren Status und ihre Modalität: Handelt es sich
beispielsweise um eine Anforderung oder eine offene Frage, ist sie
entschieden oder ungeklärt, und wird sie als verpflichtend oder als
Wunsch formuliert? Eine \emph{Disposition} kennzeichnet zusätzlich
für jedes Zielartefakt, ob der Claim dort erforderlich, als Kontext
relevant oder nicht einschlägig ist. Im Besuchsfall ist der Claim
für die Anforderungsableitung als erforderlich, für Architektur als
Kontext und für offene Fragen als nicht einschlägig eingeordnet.
Diese Zuordnungen steuern die Weiterverarbeitung; sie erteilen keine
fachliche Freigabe und garantieren keine korrekte Ableitung.

Abbildung~\ref{fig:evidenz-provenienzmodell} zeigt das logische
Modell: Durchgezogene Verbindungen beschreiben die Verarbeitung,
gestrichelte den deklarierten Herkunftsbezug. Die tatsächliche
Reihenfolge und die jeweiligen Prüf- und Klärungsbedingungen stehen
in Tabelle~\ref{t:ledger-schritte}.

\begin{figure}[htbp]
  \centering
  \includegraphics[width=\textwidth]{figures/05/Abbildung5.3.pdf}
  \caption{Logisches Evidenz- und Provenienzmodell des Evidence
    Ledgers (eigene Darstellung).}
  \label{fig:evidenz-provenienzmodell}
\end{figure}

% ============================================================
% Tabelle 5.3-1: Verarbeitungsschritte des Evidence Ledgers — v01 (09.09.)
% Label: t:ledger-schritte
% Stil: chap6.1-Konvention.
% HERKUNFT: Reihenfolge = tatsächliche Graphverdrahtung
% LedgerBuilderWorkflow.cs (Builder „LedgerBuilderUnitCoverage":
% Segmentation → Extraction → UnitCoverageGate → UnusedTriage →
% UnusedCompare → CompareRepair → Canonicalization → [Check ⇄ Repair]
% → FacetValidation); Art je Schritt = F-1/F-2 des Umsetzungsplans
% (FacetValidator = LLM-Call, det. Antwortprüfung + Statusableitung
% grounded→approved; UnusedUnitCompareRepair = det. Referenzpflicht-
% Abstufung); Queue-Filter = LedgerAdjudicationAdapter
% (DecidableCompareVerdicts); Leer-Queue-Skip = R-50
% (AdjudicationStageExecutor). Funktional gebündelt, kein
% Executor-Inventar.
% NACHZUG 09.09.: letzte Zeile = A10-Verfeinerung
% (PipelineAdjudicationRefine, geteilte Naht AdjudicationRefine.cs) —
% am 09.09. nach Abschluss der Messkampagne eingebunden
% (Standabgrenzung in 5.6; Kap.-7-Läufe entstanden DAVOR).
% ============================================================

\begin{table}[htbp]
  \centering
  \small
  \begin{tabular}{>{\raggedright\arraybackslash}p{0.235\textwidth}>{\raggedright\arraybackslash}p{0.135\textwidth}>{\raggedright\arraybackslash}p{0.51\textwidth}}
    \hline
    \textbf{Schritt (Eingabe $\rightarrow$ Ergebnis)} & \textbf{Art} & \textbf{Fortsetzung / menschliche Klärung} \\
    \hline
    Unit-Segmentierung (Transkript $\rightarrow$ adressierbare Units) & deterministisch & immer weiter \\
    Kandidaten-Extraktion (Units $\rightarrow$ quellgebundene Claim-Kandidaten) & modellgestützt & weiter \\
    Unit-Bilanz (verwendete / unverwendete Units) & deterministisch & weiter; Grundlage der Hinweisprüfung \\
    Hinweisprüfung ungenutzter Units (Triage und Vergleich gegen die Kandidaten) & modellgestützt; anschließend deterministische Belegpflicht & Prüfergebnisse je Unit (fehlender Claim · menschlich zu klären · als Beleg anzuhängen · bereits mittelbar abgedeckt); ein Deckungsurteil ohne benannten Kandidaten wird deterministisch zu „menschlich zu klären`` abgestuft \\
    Kanonisierung (Kandidaten $\rightarrow$ kanonischer Claim-Bestand, Bündelung möglich) & modellgestützt & deterministischer Übergangs-Check prüft die \emph{referenzielle} Deckung (jeder Kandidat referenziert?); bei Verstoß Reparaturschleife, sonst weiter \\
    Facetten-Validierung (Facetten und Dispositionen je Claim) & Modell-Urteil; deterministische Antwortprüfung und Statusableitung & als gestützt beurteilte Claims werden freigegeben; alle übrigen werden zur menschlichen Nachprüfung markiert \\
    Vorlagen-Aufbau (Adjudikations-Liste) & deterministisch & markierte Claims und Hinweise werden vorgelegt; „bereits mittelbar abgedeckt``-Hinweise werden nicht vorgelegt \\
    Adjudikation und Anwendung & Mensch; deterministische Anwendung & Entscheidung je Position; Ergebnis ist die verwendbare Evidenzsicht --- eine leere Vorlagenliste überspringt den Entscheidungspunkt \\
    Facetten-Vervollständigung neuer Claims & Modell-Urteil (nur bei ausstehenden Facetten) & vervollständigt Facetten und Dispositionen per Adjudikation neu aufgenommener Claims vor der fachlichen Ableitung; ohne ausstehende Facetten kein Modellaufruf (Durchreich-Schritt) \\
    \hline
  \end{tabular}
  \caption[Verarbeitungsschritte des Evidence Ledgers]{Die tatsächlich realisierten Verarbeitungsschritte des
    Evidence Ledgers in Ausführungsreihenfolge (funktional gebündelt).
    „Modellgestützt`` bezeichnet ein semantisches Modellurteil;
    deterministische Schritte prüfen Form, Referenzen oder leiten
    Status ab. Die Hinweisprüfung arbeitet auf dem
    \emph{Kandidaten}-Bestand vor der Kanonisierung.}
  \label{t:ledger-schritte}
\end{table}


Die Unit-Bilanz stellt für die Hinweisprüfung eine explizite Menge
ungenutzter Quellstellen bereit. Modelle beurteilen deren Relevanz
und Bezug zum Kandidatenbestand; der Workflow legt fest, wann diese
Urteile eingeholt und wie sie weitergegeben werden. So wird die
offene Suche nach Auslassungen für diesen Prüfkanal in eine
Bearbeitung benannter Quellstellen überführt.

% P10-6-Nachzug, 16.09.2026: Erklärabsatz nicht durch eine ganzseitige Übersicht trennen.
\begin{samepage}
Die Reihenfolge begrenzt das Blickfeld der Prüfungen: Hinweise auf
ungenutzte Quellstellen entstehen auf dem Kandidatenbestand
\emph{vor} der Kanonisierung. Inhalte, die erst bei einer späteren
Bündelung verloren gehen, werden damit nicht erneut gesucht. Der
Übergangs-Check prüft anschließend, ob Kandidaten referenziert
bleiben, nicht ob ihre Bedeutung vollständig erhalten wurde.
Abschnitt~\ref{sec:eval-stufen} untersucht diese Grenze am Lauf.
Auch die Vergabe der Dispositionen bleibt ein Modellurteil unter
expliziten Vorgaben: Dem Fragen-Pfad sollen nur tatsächlich offene
Inhalte als erforderlich zugeordnet werden. Ob etwa eine
rhetorische oder bereits geklärte Frage richtig erkannt wurde, ist
dadurch nicht gesichert.
\par
\end{samepage}

Die \emph{Evidenz-Adjudikation} ist die menschliche Entscheidung
über die vom System vorgelegten Prüf- und Klärungsfälle. Sie betrifft
die \emph{verwendbare Evidenzsicht}, also den Claim-Bestand für die
nachgelagerte fachliche Ableitung. Die Auswahl ist selektiv:
Unbeanstandete Claims werden ohne Einzelentscheidung weitergeführt;
bei leerer Vorlagenliste entfällt der Entscheidungspunkt.
Ob eine aus diesen Claims abgeleitete Änderung in den Core
übernommen werden soll, entscheidet der gesonderte
Autorisierungspfad aus Abschnitt~\ref{sec:governance-mutation}.
Ein adjudizierter Claim ist noch keine autorisierte fachliche
Zustandsänderung.

Was eine solche Vorlage enthält, zeigt die \emph{unbeantwortete}
Position AU-0124 aus der rekonstruierten Vorlagenliste des
Stresslaufs (Abschnitt~\ref{sec:eval-luecken}). Sie enthält den
Formulierungsvorschlag „Bei Auswahl eines Profils soll eine
Detailansicht erscheinen, die das Profilbild vergrößert zeigt und
die vier Hauptbereiche der App als interaktive Buttons anbietet``.
Das Systemurteil lautet \emph{fehlender Claim}; die Begründung
verweist auf eine konkrete Anforderung über die bereits erfasste
Profilübersicht hinaus. Für diesen Vorlagentyp bietet die
Bedienoberfläche fünf Aktionen: als neuen Claim übernehmen; als
Beleg an einen bestehenden Claim anhängen; einem bestehenden Claim
zuordnen (Audit); mit Begründung verwerfen; offen lassen. Diese
Optionen sind typabhängig. Welche davon hier gewählt werden sollte,
bleibt offen; eine ausgeführte Antwort wird nicht behauptet.

Nur die erste dieser Aktionen erzeugt einen zusätzlichen Claim.
Dessen Facetten und Dispositionen sind zunächst vorläufig belegt.
Im Gesamtablauf vervollständigt ein nachgelagerter modellgestützter
Schritt diese Zuordnungen vor der fachlichen Ableitung; ohne
ausstehende Facetten erfolgt kein zusätzlicher Modellaufruf.
Der Schritt garantiert keine semantisch richtige Übernahme in ein
Folgeartefakt. Die zuvor separat verfügbare Vervollständigung wurde
am 09.09.2026 im Zuge des nachträglichen Systemreviews in den
Gesamtablauf eingebunden. Der Integrationsbeleg in
Anhang~\ref{anh:evaluation} dokumentiert Einbindung und Ausführung;
die in Kapitel~\ref{chap:evaluation} ausgewerteten Ledger-Läufe
entstanden davor und belegen ihre Wirkung nicht.

Nachgelagerte Artefaktelemente übernehmen Claim-Kennungen als
Herkunftsreferenzen. Im Besuchsfall enthält ein Claim sowohl die
Ankündigung als auch die Übersicht für Pflegende. Die fachliche
Ableitung bildet daraus zwei Anforderungen, die denselben Claim
referenzieren (Tabelle~\ref{t:beispiel-besuchsplanung}). Bei
technisch auflösbaren Referenzen lässt sich so vom abgeleiteten
Element über den archivierten Claim und seine Unit-Bezüge zur
Quellstelle zurückgehen. Die deklarierte Herkunft gehört damit zur
Verarbeitung und muss nicht erst aus dem fertigen Text erschlossen
werden.

Dieser Herkunftsbezug hat klare Grenzen. Der Ledger ist der
Evidenzpfad des Transkripteingangs; andere Eingänge erhalten ihre
Herkunft durch die jeweilige kontrollierte Aufbereitung
(Abschnitt~\ref{sec:logische-gesamtarchitektur}). Eine vorhandene
Referenz garantiert weder ihre technische Auflösbarkeit noch die
semantische Stützung der Aussage oder die Erfassung aller relevanten
Inhalte. Diese Ebenen aus Abschnitt~\ref{sec:traceability-provenienz}
werden deshalb getrennt betrachtet; Kapitel~\ref{chap:evaluation}
untersucht die ausgewiesenen Eigenschaften. Die fachliche Güte der
Facetten- und Dispositionszuordnung wird nicht gesondert gemessen;
Abschnitt~\ref{sec:eval-anschluss} ordnet ihre Bedeutung für die
Weiterverarbeitung qualitativ ein.

Der Ledger macht somit semantische Zwischenstände und ihre Herkunft
überprüfbar. Die von A1 zusätzlich geforderte Beobachtung tatsächlich
ausgeführter Agenten-, Werkzeug- und Workflowaktionen ist eine
gesonderte Realisierungsaufgabe
(Abschnitt~\ref{sec:realisierung-observability}). Auf der
bereitgestellten Evidenzsicht setzt nun die fachliche Erzeugung,
Prüfung und Reparatur aus
Abschnitt~\ref{sec:saeule-semantische-transformation} auf.
